\DocumentMetadata{lang=es-DO,testphase={phase-III,math}}
\documentclass[11pt]{scrartcl}
\usepackage{icv}

% -----------------------------------------------------------------------
% Metadata -- reemplaza estos valores para cada especificación de proyecto
% nueva. "titulo" suele ser el nombre del proyecto de curso.
% -----------------------------------------------------------------------
\icvsetup{
  asignatura  = {ICV 442-T -- Acueductos y Alcantarillados},
  titulocorto = {Términos de referencia},
  titulo      = {Términos de referencia},
  subtitulo   = {Nombre del proyecto de curso},
  profesor    = {Dr. Ricardo Hernández Moreira},
  periodo     = \icvciclo{2026}{3},
  version     = {v1},
}

\begin{document}
\icvmaketitle

\begin{icvobjectives}
  \item Primera competencia o resultado de aprendizaje que evalúa el proyecto.
  \item Segunda competencia o resultado de aprendizaje.
\end{icvobjectives}

\section{Propósito del proyecto}

Descripción breve de qué se pide diseñar/producir y cuál es el producto
final esperado.

\section{Alcance}

\subsection{Componentes incluidos}
\begin{itemize}
  \item Primer componente del alcance.
  \item Segundo componente del alcance.
\end{itemize}

\subsection{Nivel de detalle esperado}
\begin{itemize}
  \item Primer criterio de nivel de detalle (p.ej. memorias de cálculo trazables).
  \item Segundo criterio.
\end{itemize}

\section{Marco normativo}
\begin{enumerate}
  \item Primera norma o reglamento aplicable.
  \item Segunda norma o reglamento aplicable.
\end{enumerate}

\begin{icvwarning}[Atención]
No basta citar una institución o norma de forma genérica -- las
decisiones de diseño deben vincularse a artículos, tablas o criterios
identificables de la norma consultada.
\end{icvwarning}

\section{Información disponible}

\begin{center}
\begin{tabular}{p{0.29\linewidth} p{0.18\linewidth} p{0.37\linewidth}}
  \toprule
  \textbf{Dato} & \textbf{Estado} & \textbf{Observación} \\
  \midrule
  Dato 1 & Disponible & Observación breve \\
  Dato 2 & Parcial     & Observación breve \\
  \bottomrule
\end{tabular}
\end{center}

\section{Organización de equipos}
\begin{itemize}
  \item Tamaño y composición de los equipos.
  \item Roles sugeridos o requeridos.
\end{itemize}

\section{Entregables}

\begin{center}
\begin{tabular}{p{0.33\linewidth} p{0.17\linewidth} p{0.30\linewidth}}
  \toprule
  \textbf{Entregable} & \textbf{Formato} & \textbf{Observaciones} \\
  \midrule
  Entregable 1 & PDF & Observación breve \\
  Entregable 2 & PDF + archivo nativo & Observación breve \\
  \bottomrule
\end{tabular}
\end{center}

\begin{icvnote}[Convención de nombres para entregas]
\texttt{CODIGO\_T[Equipo]\_[Entregable]\_[Apellidos]\_v[versión].pdf}
\end{icvnote}

\section{Criterios generales de aceptación}
\begin{itemize}
  \item Primer criterio de aceptación.
  \item Segundo criterio de aceptación.
\end{itemize}

\end{document}
